\section{Results and Conclusions}\label{sec:results}

\begin{table}[H]\centering\small
\caption{Answers to the research questions.}\label{tab:answers}
\begin{tabularx}{\linewidth}{@{}L{2.1cm}Y L{3.8cm}@{}}
\toprule
RQ & Answer & Strength of evidence\\
\midrule
RQ1 Demand \& pay & Data roles ask for a toolkit (SQL 16.5\%, SAS 13.4\%, Python 12.6\%, ML 12.0\%).
Beyond experience (OR 1.66 per year), predictive modelling, NLP, statistics, SAS and BI tools sit in
higher salary bands; MIS and generic communication mark lower-paid support roles. &
Strong: 14,840 postings, adjusted models, FDR control; salary only in bands.\\
\addlinespace
RQ2 Junior growth & Storytelling (0.59) and maths-stats (0.55) separate high- from low-hike juniors
most; model AUC 0.90 vs 0.50 baseline. & Strong but small sample (139); ceiling effect.\\
\addlinespace
RQ3 Senior success & Openness (0.77) and conscientiousness (0.76) dominate; AUC 0.96. & Strong;
unusually clean separation suggests sample data.\\
\addlinespace
RQ4 Decision & Teach Statistics, SAS, ML Foundations, Python \& SQL (robust in 9/9 scenarios) plus
Tableau storytelling (7/9); expand trainer capacity before budget. & Optimal for the stated cost and
coverage assumptions (not for real-world costs in general); robust to evidence uncertainty.\\
\bottomrule
\end{tabularx}
\end{table}

\textbf{Overall conclusion.} The evidence points in one direction from three independent angles:
\emph{quantitative foundations} (maths \& stats) are valued by the market and by employers after
hiring; \emph{coding} is the entry ticket; and \emph{communicating results} (dashboards \&
storytelling) is the skill employers reward internally but rarely advertise. For seniors, openness
and conscientiousness, rather than any single technical skill, are associated with success. An
institution can act on this immediately: a five-course plan closes most of an illustrative cohort's
gaps for Rs~3.53 lakh, and the next improvement depends on trainer capacity rather than money.

\textbf{Link to the analytics objective.} The objective (Section~\ref{sec:problem}) asked for
the skills and traits associated with demand, pay and success, and for an explainable training
plan. Each success criterion is met: every RQ has a quantified answer (Table~\ref{tab:answers});
both employee models beat the majority-class baseline under cross-validation (AUC 0.90 and 0.96 vs
0.50, permutation $p = 0.002$); and the RQ4 plan states its binding constraint (trainer-hours) and
exactly what extra resources would buy (Table~\ref{tab:hours}).

\textbf{What we did not find.} Big-data skill alone is not associated with junior hikes, and it has
the weakest combined evidence of the five areas; it enters the plan only once extra trainer capacity
is available. Meaning-based extraction from job descriptions did not work on this data and was
excluded rather than reported.
